\documentclass[book.tex]{subfiles}

\begin{document}
\chapter{Renormalization and the Renormalization Group}

\label{chap:renormalization_group}
\noindent\rule{11cm}{0.4pt} \\
\textbf{Prerequisites:} \autoref{chap:qft_basics}  \\
\textbf{Difficulty Level:} ** \\
\noindent\rule{11cm}{0.4pt}
\vspace{5mm}


Renormalization provides systematic tools for handling the infinities that arise in Feynman integrals. The renormalization group provides a structured method for looking at a physical system at multiple length scales and serves as a structural explanation of the techniques that arise in renormalization. The mathematical tools of the renormalization group come up in statistical mechanics and quantum field theory in particular.

\section{Renormalization}

There exist a number of methods to systematically cancel the infinities by the introduction of counterterms. Such counterterms are closely tied to the regularization methods that we learned about in the previous chapter, with counterterms typically dependent on the cutoff parameter $\Lambda$. In this section, we briefly survey such infinities.

\subsection{On-shell Scheme}

In the on-shell renormalization scheme, parameters in the Lagrangian are replaced with ``physical'' parameters that are set by reference to experimental data.

\subsection{Minimal Subtraction Scheme}
The minimal subtraction scheme follows a similar procedure, but only absorbs divergent terms into the counterterms.



\section{Basic Formulation of Renormalization Group}

Suppose that we have state variables $\{s_i\}$ and coupling constants $\{J_k\}$. Let us suppose the system state is then defined by a Hamiltonian
\begin{align}
    H(\{s_i\}, \{J_k\})
\end{align}
Suppose then that we have a transformation
\begin{align}
    \{s_i\} \mapsto \{\tilde{s_i}\}
\end{align}
where 
\begin{align}
    |\{s_i\}| > |\{\tilde{s}_i\}|
\end{align}
That is, the transformation shrinks the number of state variables. If we have a function $\beta$
\begin{align}
    \beta(\{Z_k\}) = \{\tilde{Z}_k\}
\end{align}
such that
\begin{align}
    H(\{s_i\}, \{Z_k\}) = H(\{\tilde{s}_i\}, \{\tilde{Z}_k\})
\end{align}
then we say that the system is renormalizable.
%\textcolor{red}{(TODO: Cite Shankar). Do I want to integrate in the Shankar Gaussian example?}

\end{document}